% !TEX program = lualatex
% ============================================================
%  Math 104 — Comprehensive Solutions
%  Ross, Elementary Analysis: The Theory of Calculus (2nd ed.)
%  Keshav Ramamurthy — Fall 2026
% ------------------------------------------------------------
%  HOW TO USE (one growing document, one banner per chapter)
%
%    Find the chapter banner below and type directly under it:
%
%        \exercise{8.6}
%        \begin{solution}
%          your write-up ...
%        \end{solution}
%
%    Restate a problem first (optional, renders a gold box):
%
%        \begin{problem}[short title]
%          statement ...
%        \end{problem}
%
%    Structure inside a solution:  claim / idea / proofsketch
%    environments, plus defn, thm, lem, prop, cor, ex, rem
%    (numbered within the current chapter).
%    Cross-reference:   \exref{8.6}
% ============================================================
\documentclass[11pt]{article}

\usepackage{fontspec}
\usepackage{microtype}
\usepackage{mathtools,amssymb,amsfonts,amsthm,bm}
\usepackage{enumitem}
\usepackage{booktabs,array,xcolor}
\usepackage{geometry,fancyhdr,setspace}
\usepackage{etoolbox}
\usepackage[most]{tcolorbox}
\usepackage[hidelinks]{hyperref}
\usepackage[nameinlink,noabbrev]{cleveref}

% ---------- Page and paragraph rhythm (matches lecture template) ----------
\geometry{margin=1in,headheight=15pt}
\setstretch{1.12}
\setlength{\parindent}{0pt}
\setlength{\parskip}{0.55em plus 0.12em minus 0.08em}
\setlength{\abovedisplayskip}{0.8em plus 0.2em minus 0.1em}
\setlength{\belowdisplayskip}{0.8em plus 0.2em minus 0.1em}
\allowdisplaybreaks[3]
\setlist{leftmargin=*,topsep=0.35em,itemsep=0.15em,parsep=0pt}

% ---------- Palette (matches reference cheatsheet) ----------
\definecolor{ink}{HTML}{1E293B}
\definecolor{navy}{HTML}{183B56}
\definecolor{paperblue}{HTML}{E8F1F8}
\definecolor{gold}{HTML}{FFF3CD}
\definecolor{mint}{HTML}{E8F5E9}
\definecolor{muted}{HTML}{52606D}
\definecolor{rule}{HTML}{C9D5DF}
\color{ink}

% ---------- Metadata ----------
\newcommand{\studentname}{Keshav Ramamurthy}
\newcommand{\coursename}{Math 104}
\newcommand{\courselongname}{Real Analysis}
\newcommand{\textbookname}{Ross, \emph{Elementary Analysis: The Theory of Calculus} (2nd ed.)}

% ---------- Core notation (same as ~/.config/nvim/textemplate.tex) ----------
\newcommand{\N}{\mathbb{N}}
\newcommand{\Z}{\mathbb{Z}}
\newcommand{\Q}{\mathbb{Q}}
\newcommand{\R}{\mathbb{R}}
\newcommand{\C}{\mathbb{C}}
\newcommand{\F}{\mathbb{F}}
\newcommand{\K}{\mathbb{K}}
\newcommand{\E}{\mathbb{E}}
\newcommand{\Prob}{\mathbb{P}}
\newcommand{\1}{\mathbf{1}}
\newcommand{\eps}{\varepsilon}
\newcommand{\dd}{\,\mathrm{d}}
\newcommand{\abs}[1]{\left\lvert#1\right\rvert}
\newcommand{\norm}[1]{\left\lVert#1\right\rVert}
\newcommand{\ip}[2]{\left\langle#1,#2\right\rangle}
\newcommand{\set}[1]{\left\{#1\right\}}
\newcommand{\ceil}[1]{\left\lceil#1\right\rceil}
\newcommand{\floor}[1]{\left\lfloor#1\right\rfloor}
\newcommand{\given}{\,\middle\vert\,}
\newcommand{\indep}{\perp\!\!\!\perp}
\newcommand{\wh}[1]{\widehat{#1}}
\newcommand{\deriv}[2]{\frac{\mathrm{d}#1}{\mathrm{d}#2}}
\newcommand{\pderiv}[2]{\frac{\partial#1}{\partial#2}}
\newcommand{\lto}{\longrightarrow}
\newcommand{\st}{\text{ such that }}
\DeclareMathOperator{\Aut}{Aut}
\DeclareMathOperator{\End}{End}
\DeclareMathOperator{\Gal}{Gal}
\DeclareMathOperator{\Hom}{Hom}
\DeclareMathOperator{\Ker}{ker}
\DeclareMathOperator{\im}{im}
\DeclareMathOperator{\rank}{rank}
\DeclareMathOperator{\nullity}{nullity}
\DeclareMathOperator{\Span}{span}
\DeclareMathOperator{\tr}{tr}
\DeclareMathOperator{\diag}{diag}
\DeclareMathOperator{\spec}{spec}
\DeclareMathOperator{\ord}{ord}
\DeclareMathOperator{\supp}{supp}
\DeclareMathOperator{\diam}{diam}
\DeclareMathOperator{\Var}{Var}
\DeclareMathOperator{\Cov}{Cov}
\DeclareMathOperator*{\argmin}{arg\,min}
\DeclareMathOperator*{\argmax}{arg\,max}

% ---------- Course-specific notation (analysis) ----------
\newcommand{\limsupn}[1]{\limsup_{n\to\infty} #1}
\newcommand{\liminfn}[1]{\liminf_{n\to\infty} #1}
\newcommand{\ball}[2]{B_{#2}(#1)}                 % ball center #1, radius #2
\newcommand{\eqdef}{\coloneqq}

% ---------- Theorem-like environments (numbered by chapter) ----------
\newcounter{chap}
\newtheoremstyle{notestyle}{0.7em}{0.7em}{\itshape}{}{}{\bfseries}{.5em}{}
\theoremstyle{notestyle}
\newtheorem{theorem}{Theorem}[chap]
\newtheorem{lemma}[theorem]{Lemma}
\newtheorem{proposition}[theorem]{Proposition}
\newtheorem{corollary}[theorem]{Corollary}
\theoremstyle{definition}
\newtheorem{definition}[theorem]{Definition}
\newtheorem{example}[theorem]{Example}
\newtheorem{remark}[theorem]{Remark}

% ---------- Problem machinery ----------
% Chapter banner: \chapterbanner{2}{Sequences}{Ross, \S\S 7--16}
\newtcolorbox{chapterbox}[2]{%
  enhanced,
  colback=paperblue, colframe=navy, boxrule=0.7pt, arc=2.5pt,
  left=10pt, right=10pt, top=12pt, bottom=8pt,
  boxed title style={colback=navy, colframe=navy, arc=2pt, boxrule=0pt,
    left=8pt, right=8pt, top=3pt, bottom=3pt},
  attach boxed title to top left={yshift=-2.4mm, xshift=5mm},
  title={\sffamily\bfseries\color{white}\large Chapter #1\quad #2}}
\newcommand{\chapterbanner}[3]{%
  \clearpage
  \setcounter{chap}{#1}%
  \setcounter{theorem}{0}%
  \phantomsection
  \addcontentsline{toc}{section}{Chapter #1\quad #2}%
  \begin{chapterbox}{#1}{#2}
    {\small\sffamily\color{muted}#3\par}
  \end{chapterbox}
  \medskip\nopagebreak}

% Exercise header: \exercise{8.6}  →  bold header + thin rule, then type
\newcommand{\exercise}[1]{%
  \par\goodbreak\medskip\noindent
  {\sffamily\bfseries\color{navy}Exercise #1}%
  \par\nobreak\vspace{1.5pt}
  {\color{rule}\hrule height 0.7pt}\par\nobreak\smallskip\nopagebreak}

% Cross-reference an exercise: \exref{8.6}
\newcommand{\exref}[1]{\textsf{\textbf{Exercise~#1}}}

% Problem statement box: \begin{problem}[short title] ... \end{problem}
\newtcolorbox{problem}[1][]{%
  enhanced, breakable,
  colback=gold!50, colframe=rule, boxrule=0.5pt, arc=2pt,
  colbacktitle=gold!50, coltitle=navy,
  left=8pt, right=8pt, top=6pt, bottom=6pt,
  before skip=8pt, after skip=8pt,
  title={\sffamily\bfseries Problem\ifstrempty{#1}{}{\ ---\ #1}}}

% Solution / answer blocks (qed square at the end)
\newenvironment{solution}{\begin{proof}[Solution]}{\end{proof}}
\newenvironment{answer}{\begin{proof}[Answer]}{\end{proof}}
\renewcommand{\qedsymbol}{$\square$}

% Quick structure inside solutions
\newenvironment{claim}{\par\smallskip\noindent\textit{Claim.}\ }{\par\smallskip}
\newenvironment{idea}{\par\smallskip\noindent\textit{Idea.}\ }{\par\smallskip}
\newenvironment{proofsketch}{\par\smallskip\noindent\textit{Proof sketch.}\ }{\par\smallskip}

% ---------- Header ----------
\pagestyle{fancy}
\fancyhf{}
\fancyhead[L]{\coursename}
\fancyhead[C]{\courselongname\ — solutions}
\fancyhead[R]{\studentname}
\fancyfoot[C]{\thepage}

\begin{document}

% ---------- Title ----------
\begin{center}
  {\Large\sffamily\bfseries\color{navy}\coursename\ -- \courselongname}\\[3pt]
  {\sffamily\color{muted} Comprehensive exercise solutions}\\[3pt]
  {\small\sffamily\color{muted}\textbookname\quad\textbullet\quad Fall 2026\quad\textbullet\quad \studentname}
\end{center}
\medskip
{\color{rule}\hrule height 0.8pt}
\bigskip

{\small
\begin{tcolorbox}[colback=paperblue, colframe=rule, boxrule=0.4pt, arc=2pt,
  left=7pt, right=7pt, top=5pt, bottom=5pt,
  title={\sffamily\bfseries\color{navy} Adding a solution},
  colbacktitle=paperblue]
\sffamily To record a solved exercise, type under its chapter banner:
\begin{center}
\texttt{\textbackslash exercise\{8.6\}\ \ \textbackslash begin\{solution\}\ ...\ \textbackslash end\{solution\}}
\end{center}
Restate the problem first with \texttt{\textbackslash begin\{problem\}[title]...\textbackslash end\{problem\}};
structure proofs with \texttt{claim}, \texttt{idea}, \texttt{proofsketch}.
\end{tcolorbox}
}

\medskip
\tableofcontents

% =====================================================================
%  CHAPTER 1 — INTRODUCTION
% =====================================================================
\chapterbanner{1}{Introduction}{Ross, \S\S 1--6:\ natural, rational, and real numbers; the completeness axiom}
\begin{example}[1.1]
We proceed via induction.
Let $a_n = \sum_{i=1}^{n} i^2$.\par  Then $a_1 = 1 = \frac{1}{6}(1)(2)(3)$, so the base case
    holds.
    \par Then, we proceed via induction. Suppose $a_i = \frac{1}{6}i(i + 1) (2i + 1)$
    for $i \in \{1, 2, 3, 4, \cdots, k\}.$ 
    \par We will show that this implies that $a_{k +
    1} = \frac{1}{6}(k + 1)(k + 2)(2k + 3)$
    \par We add $a_{k + 1} = a_{k} + (k + 1)^2 = \frac{1}{6} (k)(k+1)(2k+1) + (k+1)^2$
    This is from our inductive hypothesis, and yields $$a_{k+1} = (k +
    1)\left(\frac{1}{6}(k)(2k + 1) + k + 1\right)$$
    Pulling out $\frac{1}{6}$ yields $\frac{1}{6}(k + 1)(k(2k + 1) + 6k + 6)$
    And the second term of that polynomial is $$2k^2 + 7k + 6 = 2k^2 + 4k + 3k + 6 = (2k
    +3)(k + 2)$$
    satisfying our inductive hypothesis.
\end{example}
% REST OF CHAPTER ONE IS EASY, SO CONTINUING TO CHAPTER 2 FOR 104

% \exercise{1.1}
% \begin{solution}
%   ...
% \end{solution}

% =====================================================================
%  CHAPTER 2 — SEQUENCES
% =====================================================================
\chapterbanner{2}{Sequences}{Ross, \S\S 7--16:\ limits of sequences, limit theorems, monotone/Cauchy, subsequences, limsup/liminf, series}
\begin{example}[2.2]
Show that $\sqrt[3]{2}, \sqrt[7]{5}, \sqrt[4]{13}$ are not rational.
\par For the first for sake of contradiction set $\sqrt[3]{2} = \frac{p}{q}$ for $p, q \in \mathbb{Z}^{+}, \gcd(p, q)
= 1.$
Then, cubing both sides, we have $2 = \frac{p^3}{q^3} \rightarrow p^3 = 2q^3.$
\par However, since the rhs is divisible by 2, we have $2 \mid p \rightarrow p = 2p', p'
\in \Z^{+}.$
\par Substituting this in yields $$(2p')^3 = 2q^3 \rightarrow 4(p')^3 = q^3$$
So the analogous idea implies $2 \mid q$. However this is a contradiction, as $\gcd(p,
q) = 1$ but we found 2 divides both $p$ and $q$.

\par We can show the others analogously.
\end{example}
\begin{example}[2.3]
Show that $\sqrt{2 + \sqrt2}$ is not rational.
Suppose for sake of contradiction we set $\sqrt{2 + \sqrt{2}} = \frac{p}{q}$ again, we assume $p, q \in \Z^{+}$ and $\gcd(p, q) = 1$.
Squaring, we immediately have a contradiction because $$2 + \sqrt2 = \frac{p^2}{q^2}
\rightarrow \sqrt{2} = \frac{p^2}{q^2} - 2$$
and the lhs is irrational but the rhs is rational.
\end{example}
% \exercise{7.1}
% \begin{solution}
%   ...
% \end{solution}

% =====================================================================
%  CHAPTER 3 — CONTINUITY
% =====================================================================
\chapterbanner{3}{Continuity}{Ross, \S\S 17--22:\ continuous functions, properties, uniform continuity, limits of functions}

\begin{example}[3.1]
For the set $\N$, we have: (for arbitrary $a, b, c \in \N$)
\begin{itemize}
    \item a1: associativity holds $(a + b) + c = a+(b+c)$
    \item a2: commutativity holds. $a + b = b + a$
    \item a3: identity exists (0)

\end{itemize}

We continue and we see the existence of a multiplicative inverse, namely for
$a \in \N$, $\exists a^{-1} \text{ s.t. } aa^{-1} = 1$ does not hold. For example
consider $a = 2$ and $a^{-1} = \frac{1}{2} \notin \N$.
$\Z$ similarly does not have a multiplicative inverse for all $a \in \Z$.
\end{example}
\begin{example}[3.2]
Commutativity of addition is shown in the proof $(vi)$ to show $b^{-1} * ab = ab *
b^{-1} = a$
\par It's shown in part $(ii)$, while adding $0$ at the beginning and then while
adding to the rhs adding it to the right side (which is commutativity) then cancelling
the zero out.
\par Shown in $(iii)$, where $0 \cdot b = b \cdot 0 = 0.$
\end{example}
\begin{example}[3.3]
We wish to show that $ac= bc$ and $c \ne 0$ implies $a = b.$
We have $$ac + (-ac) = bc + (-ac) \rightarrow 0 = c(b - a)$$
Then we write $$c^{-1} \cdot 0 = b - a \rightarrow 0 = b - a$$
Adding $a$ on both sides yields $a = b$.

 To show $ab = 0$ implies $a = 0$ or $b = 0$ for $a, b \in \R$
\par We proceed by contradiction. Suppose $a \ne 0, b \ne 0.$
Notice we can write $$a^{-1} (ab) b^{-1} = a^{-1} \cdot 0 \cdot b^{-1}$$
So the lhs is 1 but the rhs is 0, a contradiction.
\end{example}
\begin{example}[3.4]
To show $0 < 1$ we can suppose for sake of contradiction $0 \ge 1$
Then, we'd have $-1 \ge 0$, adding $-1$ to both sides
Multiplying both sides by $(-1)$ we'd have $$(-1)(-1) \le 0 \rightarrow 1 \le 0.$$
Even though we used less than or equal to, obviously $0 \ne 1$ by definition, so we have
a contradiction.
Given the ordering properties, we wish to show if $0 < a$, then $0 < a^{-1}.$
% We have $a a^{-1} = 1 \rightarrow a^{-1} = \frac{1}{a}.$
\par Suppose for sake of contradiction $a^{-1} \le 0.$ Then we have, multiplying a on both
sides $$a \cdot a^{-1} \le a \cdot 0 \rightarrow 1 \le 0$$ which is a contradiction.
\par To prove $0 < a < b$ implies $0 < b^{-1} < a^{-1}$, we go in multiple steps.
\par We know that $a^{-1}$ and $b^{-1}$ are positive, we just wish to establish order.
We multiply $0 < a < b$ by the constant $a^{-1}$ to get $0 < a a^{-1} < ba^{-1}.$
\par Then multiply by $b^{-1}$ to it to get $$0 < b^{-1} a a^{-1} < b a^{-1}b^{-1}
\rightarrow 0 < b^{-1} < a^{-1}$$
% \par To show that $0 
\end{example}
% =====================================================================
%
%  CHAPTER 4 — SEQUENCES AND SERIES OF FUNCTIONS
% =====================================================================
\chapterbanner{4}{Sequences and Series of Functions}{Ross, \S\S 23--27:\ power series, uniform convergence, term-by-term calculus}

% =====================================================================
%  CHAPTER 5 — DIFFERENTIATION
% =====================================================================
\chapterbanner{5}{Differentiation}{Ross, \S\S 28--31:\ derivative basics, Mean Value Theorem, L'Hospital, Taylor's Theorem}

% =====================================================================
%  CHAPTER 6 — INTEGRATION
% =====================================================================
\chapterbanner{6}{Integration}{Ross, \S\S 32--36:\ Riemann integral, its properties, Fundamental Theorem of Calculus}

% =====================================================================
%  CHAPTER 7 — CAPSTONE
% =====================================================================
\chapterbanner{7}{Capstone}{Ross, \S\S 37--38:\ exponents and logarithms; continuous nowhere-differentiable functions}

\end{document}
